%%% Chapter 3 — Closing summary
%%% UF-F10: third chapter, satisfies the "closing summary chapter" rule.

\chapter{CONCLUSION}\label{ch:conclusion}

This dissertation presented a two-layer validator for the University of
Florida doctoral dissertation formatting specification.
Chapter~\ref{ch:intro} motivated the work and outlined contributions.
Chapter~\ref{ch:methods} described the source-layer and PDF-layer
checks, formalized the override-detection problem in
Lemma~\ref{lem:override}, stated soundness in
Theorem~\ref{thm:source-soundness}, gave the margin acceptance
criterion in Equation~\ref{eq:margin}, summarized severity tiers in
Table~\ref{tab:severity}, sketched the architecture in
Figure~\ref{fig:architecture}, and showed the input-mode routing in
Figure~\ref{fig:input-modes}. Supplementary rule identifiers are
listed in Appendix~\ref{app:rules}; synthetic broken inputs used to
verify the validator's rule coverage are described in
Appendix~\ref{app:fixtures}.

\section{Findings}

Most UF formatting rules are enforced by the \texttt{ufdissertation}
class, which means the validator's primary job is to detect overrides
of the template's correct defaults rather than to re-implement the
formatting rules themselves. A small set of rules genuinely require
inspection of the rendered PDF; these are concentrated in the
geometric and accessibility categories.

\section{Limitations}

The validator targets the Fall 2025+ UF template only. Older templates
(accepted through Summer 2026) are detected and the validator exits
with a clear message rather than producing potentially misleading
output. Master's theses are out of scope for the initial release.

\section{Future Work}

Future releases will expand input surfaces: a Chrome extension for
students writing in Overleaf, and a Visual Studio Code extension for
students using a git-plus-editor workflow. External URL liveness
checks, deferred from the initial release to keep the validator
network-free, will be added behind an opt-in flag.

\section{A Final Word}

To the student who used this demo as a reference for their own
dissertation: congratulations. The fact that you reached this section
of this document means you are within reach of submission. Take a
breath, work through the checklist in Appendix~\ref{app:checklist},
and submit when ready. The graduate school has seen better
dissertations than yours and worse, and the difference between them
was never the margins.
